% TOP_PROBLEM: {"id": 1103, "title": "Köthe Conjecture", "category_id": 4, "category": {"id": 4, "name": "algebra", "display_name": "Algebra", "description": "Group theory, ring theory, field theory, and algebraic structures.", "slug": "algebra", "order_index": 4, "created_at": "2026-07-31T15:26:25.671Z"}, "status": "open"}
% FIRST_POSTED: 2026-09-25
% ENTRY_KIND: statement_only
% !TeX program = lualatex
% Definition and sources only; this file is not an LLM solution attempt.
\documentclass[11pt]{article}
\usepackage[margin=25mm]{geometry}
\usepackage{fontspec}
\setmainfont{Latin Modern Roman}
\usepackage{amsmath,amssymb,mathtools}
\usepackage{unicode-math}
\setmathfont{Latin Modern Math}
\usepackage{xurl,hyperref}
\hypersetup{colorlinks=true,urlcolor=blue}
\setcounter{secnumdepth}{0}
\setlength{\parindent}{0pt}
\setlength{\parskip}{5pt}
\setlength{\emergencystretch}{3em}
\begin{document}
\section*{\texorpdfstring{Köthe Conjecture}{Köthe Conjecture}}\label{1103}
\subsection{Definitions and mathematical statement}
An associative ring or ideal is nil if each of its elements is nilpotent, with a nilpotence exponent that may depend on the element. This is weaker than being nilpotent, which requires one uniform bound on lengths of arbitrary products. Rings in the nil-ring formulations need not have an identity.

K\"othe's assertion is that a ring with no nonzero nil two-sided ideal also has no nonzero nil one-sided ideal. The equivalent formulations recorded in the source include: the sum of two nil right ideals is nil; and
\[
 R\text{ nil}\quad\Longrightarrow\quad M_2(R)\text{ nil}.
\]
The conclusion in the matrix version concerns every matrix individually, not a uniform common nilpotence exponent. Commutativity and finite dimensionality are not hypotheses.
\par\smallskip\textit{Formulation and concept references:} \href{https://www.mathunion.org/fileadmin/ICM/Proceedings/ICM2006.2/ICM2006.2.ocr.pdf}{[S1]}.

\subsection{Short English statement}
Can nil behavior in a one-sided ideal fail to produce a nonzero nil two-sided ideal? Equivalently, do two-by-two matrices over nil rings remain nil?
\subsection{Sources}
\begin{itemize}
\item[S1]Agata Smoktunowicz, "Some results in noncommutative ring theory," Proceedings of the ICM, Madrid 2006, Vol. II, EMS, pp. 259-269.
\url{https://www.mathunion.org/fileadmin/ICM/Proceedings/ICM2006.2/ICM2006.2.ocr.pdf}.
\textit{Formulation source.}
\end{itemize}
\end{document}
